\chapter{Curve shortening and ramp approximations}
\label{ch:ramps-v4}

\section{The moving-metric problem and its interfaces}
\label{sec:ramps-purpose}

The width argument needs a family of loops that can be evolved by curve
shortening while a Ricci metric changes.  This chapter supplies that family.
The parameter on a loop is kept fixed; it is not silently replaced by
arc-length, and the same convention is used at self-intersections.  The
ambient flow is supplied by the preceding chapters, while the annular
comparison chapter consumes the ramps constructed here.

Fix a compact Hausdorff second-countable smooth manifold $M$ of dimension
$n$, a closed interval $I=[a,b]$ with $a<b$, and an actual Ricci flow
$F=(g(t),D(t))$ on $I$.  Thus
\[
  \partial_tg(t)=-2\operatorname{Ric}_{g(t)},
\]
with $D(t)$ the Levi--Civita connection.  A curve is a displayed map
$c:\mathbb R\times I\to M$, periodic with the fixed parameter period
\[
  \mathsf p=2\pi.
\]
For a spatial derivative write $X=\partial_xc$, $v=|X|_{g(t)}$,
$S=X/v$, and use $D_s=v^{-1}D_x$.  The curvature vector and its scalar norm
are
\[
 H=D_sS,\qquad k^2=g(t)(H,H),\qquad k=\sqrt{k^2}.
\]
These are pullback-bundle quantities: two parameter values mapping to the
same point retain their own tangent vectors.  For a scalar $f(x)$ put
$\partial_sf=v^{-1}\partial_xf$, and define
\[
 L(t)=\int_0^{\mathsf p}v(x,t)\,dx,\qquad
 \Theta(t)=\int_0^{\mathsf p}k(x,t)v(x,t)\,dx.
\]
The total curvature is therefore an integral in the fixed label, not an
integral over an implicitly reparametrized image.

The two principal outputs are: (i) solutions with fixed labels for arbitrary
immersed $C^2$ data, and (ii) positive-slope graph solutions in a product
$M\times \mathbb R/(\rho\mathbb Z)$.  The latter are the ramps used to
replace a rough loop family by a smooth family in the next chapter.

\section{Corrected curve evolution}
\label{sec:ramps-corrected-evolution}

\begin{definition}[Admissible shrinking curve]
\label{def:ramps-shrinking-curve}
\lean{PoincareMT.M62ShrinkingCurve}
An admissible shrinking curve for $F$ is a periodic map $c$ such that every
slice is immersed and $C^2$, the map is smooth on
$\mathbb R\times(a,b)$, and the value, velocity, and curvature vector extend
continuously to $\mathbb R\times I$.  On the open cylinder it satisfies
\[
  \partial_tc(x,t)=H(x,t).
\]
Only the spatial $C^2$ regularity is required at the two time endpoints.
\end{definition}

For nonsmooth initial labels we use a second class: a $C^2$ shrinking curve
has the same periodicity, immersion, endpoint continuity and equation, but
is only jointly $C^1$ in the open cylinder and $C^2$ on each spatial slice.
Intrinsic positive-time regularity means that every $D_s^iH$ exists, is
jointly $C^1$ in the interior, and is continuous on each closed slab strictly
after the initial time.  It does not assert smoothness in the original
label $x$.  Composing a smooth solution with a fixed nonsmooth $C^2$
orientation-preserving parametrization explains why that distinction is
necessary.

Write $Q=\operatorname{Ric}_{g(t)}$.  Our curvature convention is
$\operatorname{Rm}(A,B,C,D)=g(R(A,B)D,C)$, where
$R(A,B)=D_AD_B-D_BD_A-D_{[A,B]}$.
The speed calculation uses $D_tX=D_xH$ and
$g(D_sH,S)=-k^2$: differentiating $g(X,X)$ gives
$-2Q(X,X)+2g(D_xH,X)=-2(Q(S,S)+k^2)v^2$.

The elementary identities already contain the moving-metric correction.  On
the interior of $I$,
\begin{equation}\label{eq:ramps-speed}
 \partial_t(v^2)=-2\bigl(\operatorname{Ric}(S,S)+k^2\bigr)v^2,
 \qquad
 \partial_tv=-\bigl(\operatorname{Ric}(S,S)+k^2\bigr)v,
\end{equation}
and hence
\begin{equation}\label{eq:ramps-commutator}
 [\partial_t,\partial_s]
   =\bigl(k^2+\operatorname{Ric}(S,S)\bigr)\partial_s
\end{equation}
as an identity on every smooth scalar test function.  The bracket is an
identity of pullback derivatives, not a bracket of vector fields on the image
of a possibly self-intersecting curve.

For the corrected curvature computation, let $P_g$ be the $g(t)$-normal
projection of $D_sH$ perpendicular to $S$.  The exact spatial equation is
\begin{equation}\label{eq:ramps-curvature}
\begin{split}
 \partial_t k^2={}&\partial_s^2k^2-2|P_g|^2+2(k^2)^2
 -2\operatorname{Ric}(H,H)+4k^2\operatorname{Ric}(S,S)\\
 &+2\operatorname{Rm}(H,S,H,S)
 +2(D_H\operatorname{Ric})(S,S)-4(D_S\operatorname{Ric})(S,H).
\end{split}
\end{equation}
Every tensor is evaluated at $(c(x,t),t)$.  Here is the calculation, including
the connection variation.  If $C(A,B)=(\partial_tD)(A,B)$, differentiation
of the Koszul formula gives
\[
 g(C(A,B),Z)=-(D_AQ)(B,Z)-(D_BQ)(A,Z)+(D_ZQ)(A,B).
\]
Put $q=k^2+Q(S,S)$.  Commutation of pullback covariant derivatives gives
\[
 D_tS=D_sH+qS,\qquad
 D_tH=D_s^2H+(\partial_sq)S+2qH+R(H,S)S+C(S,S).
\]
Differentiate $g(H,H)$, use $g(H,S)=0$, and substitute
$|D_sH|^2=|P_g|^2+k^4$ in
$\partial_s^2k^2=2g(D_s^2H,H)+2|D_sH|^2$.
This gives (\ref{eq:ramps-curvature}), with both covariant derivatives of
$Q$ present.

For comparison with the spacetime computation, equip the open manifold
$M\times(a,b)$ with $\widehat g=g(t)+dt^2$, and put $T=\partial_t$.
Horizontal lifts are understood in the following identities:
\begin{align*}
 \widehat D_AB&=D_AB+Q(A,B)T,&
 \widehat D_AT&=-Q^\sharp A,\\
 \widehat D_TB&=\partial_tB-Q^\sharp B,&
 \widehat D_TT&=0.
\end{align*}
These follow directly from the Koszul formula.  Substitution in the curvature
definition gives
\begin{align*}
 \widehat{\operatorname{Rm}}(A,B,C,D)
 &=\operatorname{Rm}(A,B,C,D)-Q(B,D)Q(A,C)+Q(A,D)Q(B,C),\\
 \widehat{\operatorname{Rm}}(T,B,C,D)
 &=(D_CQ)(B,D)-(D_DQ)(B,C).
\end{align*}
Let $\widehat H=(H,0)$ and $V=(H,1)$; these are different vectors.
If $\widehat P$ is the projection of $\widehat D_s\widehat H$
perpendicular to $S$, then
$\widehat P=P_g+Q(S,H)T$.  The full equation is
\begin{equation}\label{eq:ramps-spacetime-curvature}
\begin{split}
 \partial_tk^2={}&\partial_s^2k^2-2|\widehat P|^2+2k^4
 -2Q(H,H)+4k^2Q(S,S)\\
 &+2\widehat{\operatorname{Rm}}(V,S,\widehat H,S)
 +2Q(S,S)Q(H,H)-2(D_SQ)(S,H).
\end{split}
\end{equation}
Indeed $|\widehat P|^2=|P_g|^2+Q(S,H)^2$; the Gauss term contributes
$-2Q(S,S)Q(H,H)+2Q(S,H)^2$, and the mixed curvature contributes
$2(D_HQ)(S,S)-2(D_SQ)(S,H)$.  The quadratic Ricci terms cancel,
recovering (\ref{eq:ramps-curvature}).  The mixed curvature has a term
linear in $H$, so an error estimate proportional only to $k^2$ is
insufficient.  These are the corrections in
\cite[Lemma 0.2, pp.~4--6]{MT2015Correction}.

\begin{proposition}[Compact ambient bounds and zero-safe regularization]
\label{prop:ramps-regularization}
\uses{def:ramps-shrinking-curve}
\lean{PoincareMT.m62CorrectedCurveEvolution_from_predecessors}
\source{MT2015Correction}
\dcref{Lemma 0.2, Corollary 0.3, equation (0.4), pp. 4--7}
There are nonnegative constants $K_0,K_1,K_2$, chosen from $F$ and $I$
before choosing a curve, a regularization parameter, or a circle
circumference, such that $K_0,K_1,K_2$ bound respectively the full curvature,
$D\operatorname{Ric}$, and $\operatorname{Ric}$ on unit vectors.  Put
\[
 C_0=2K_0+6K_2+2K_2^2+6K_1,\qquad C_1=C_0/2.
\]
Then every admissible curve satisfies, on $a<t<b$,
\[
 \partial_tk^2\leq\partial_s^2k^2-2|P_g|^2+2(k^2)^2+C_0(k^2+k).
\]
For $\varepsilon>0$ define
\[
 h_\varepsilon=\sqrt{k^2+\varepsilon^2}.
\]
It is positive and smooth on the interior, obeys
\[
 (\partial_sh_\varepsilon)^2\leq|P_g|^2,\qquad
 \partial_th_\varepsilon\leq\partial_s^2h_\varepsilon+k^3+C_1(h_\varepsilon+1),
\]
and the integrated quantities satisfy
\[
 0\leq\Theta_\varepsilon(t)-\Theta(t)\leq\varepsilon L(t).
\]
\end{proposition}
\begin{proof}
Take maxima of the tensor norms on the compact unit-vector bundles over
$M\times I$.  The last five terms of (\ref{eq:ramps-curvature}) have
absolute sum at most $(2K_0+6K_2)k^2+6K_1k$.
To retain $-2|\widehat P|^2$ instead of $-2|P_g|^2$ costs at most
$2K_2^2k^2$.  Thus the stated $C_0$ works with either normal derivative.
Since $\langle H,D_sH\rangle=\langle H,P_g\rangle$,
$h_{\varepsilon,s}=\langle H,P_g\rangle/h_\varepsilon$, whence
$h_{\varepsilon,s}^2\leq|P_g|^2$.
The chain rule now yields
\[
 h_{\varepsilon,t}-h_{\varepsilon,ss}
 \leq\frac{-|P_g|^2+h_{\varepsilon,s}^2+k^4}{h_\varepsilon}
       +C_1\frac{k^2+k}{h_\varepsilon}
 \leq k^3+C_1(h_\varepsilon+1).
\]
All denominators are positive even where $H=0$.  Finally define
$\Theta_\varepsilon=\int h_\varepsilon\,ds$ and integrate
$0\leq h_\varepsilon-k\leq\varepsilon$.
See \cite[Corollary 0.3 and (0.4), pp.~6--7]{MT2015Correction}.
\end{proof}

\begin{proposition}[Length and total-curvature control]
\label{prop:ramps-length-curvature}
\lean{PoincareMT.m62CorrectedCurveEvolution_from_predecessors}
\uses{prop:ramps-regularization}
For $a\leq s\leq t\leq b$, $L$ and $\Theta$ are continuous and
\[
 L'(t)=-\int_0^{\mathsf p}(k^2+\operatorname{Ric}(S,S))v\,dx,
 \qquad L'(t)\leq K_2L(t)
\]
at interior times.  Moreover
\begin{equation}
 \Theta(t)-\Theta(s)\leq
 \int_s^t\bigl((C_1+K_2)\Theta(r)+C_1L(r)\bigr)\,dr,
 \label{eq:ramps-theta-integral}
\end{equation}
and, for $t<b$, the upper right difference quotient has the same bound up
to an arbitrary positive error.  In particular, Gronwall gives exponential
length and total-curvature bounds from the two initial quantities.
\end{proposition}
\begin{proof}
The speed identity (\ref{eq:ramps-speed}) gives the length derivative after
periodic integration.  For $\Theta_\varepsilon$, differentiate its measure
as well as its integrand:
\[
 \Theta_\varepsilon'
 \leq\int\bigl(h_{\varepsilon,ss}+k^3-h_\varepsilon k^2
          +C_1(h_\varepsilon+1)-Q(S,S)h_\varepsilon\bigr)\,ds.
\]
The cubic contribution is nonpositive.  Periodicity gives
$\int h_{\varepsilon,ss}\,ds=0$ even though $ds=v\,dx$ varies in time.
The resulting linear inequality is
\[
 \Theta_\varepsilon'\leq(C_1+K_2)\Theta_\varepsilon+C_1L.
\]
Integrate it on $[s,t]$, then use the uniform regularization error from
Proposition~\ref{prop:ramps-regularization}.  This produces
(\ref{eq:ramps-theta-integral}) without claiming that $\Theta$ is classically
differentiable where curvature vanishes.  Endpoint continuity extends the
integral inequality to $s=a$ and $t=b$.  Explicitly, for every $t<b$ and
$\eta>0$ there is $d>0$ such that, for $t<r\leq b$ and $r-t<d$,
\[
 \frac{\Theta(r)-\Theta(t)}{r-t}
 \leq(C_1+K_2)\Theta(t)+C_1L(t)+\eta.
\]
Adding the length inequality gives
\[
 L(t)\leq L(s)e^{K_2(t-s)},\qquad
 (\Theta+L)(t)\leq(\Theta+L)(s)e^{(C_1+K_2)(t-s)}.
\]
In particular the total-curvature bound depends on initial length as well
as initial total curvature, as in
\cite[Lemma 0.4, pp.~7--8]{MT2015Correction}.
The same length equation controls the time integral of curvature energy.
Put $\mathcal E(t)=\int k^2\,ds\ge0$.  Continuity of the curvature
and speed up to the endpoints makes this integrable, and
$L'(t)+\mathcal E(t)\le K_2L(t)$ implies
\[
 \frac d{dt}\bigl(e^{-K_2(t-a)}L(t)\bigr)
       \le-e^{-K_2(t-a)}\mathcal E(t).
\]
Integration and nonnegativity of the final length give
\[
 \int_a^b\mathcal E(t)\,dt
 \le e^{K_2(b-a)}
       \int_a^b e^{-K_2(t-a)}\mathcal E(t)\,dt
 \le L(a)e^{K_2(b-a)}.
\]
The endpoint passage follows by continuity from interior intervals.
For the canonical family, $L(a)$ is bounded by the original
family's length supremum plus one, uniformly for $0<\rho<1$.
This is the integral estimate used to select energy-good times.
\end{proof}

\section{The actual product circle and the slope equation}
\label{sec:ramps-circle}

For every $\rho>0$ let $C_\rho=\mathbb R/(\rho\mathbb Z)$ with its quotient
circle metric $da^2$ and its positive unit field $U$.  The product flow is the
actual product carrier $M\times C_\rho$ with metric $g(t)+da^2$ and the
product Levi--Civita connection.  The quotient construction matters: the
circle is not a coordinate interval with its endpoints identified by hand.
The product curvature, Ricci tensor, and covariant Ricci derivative are the
pullbacks of the corresponding tensors on $M$; $U$ is parallel and the same
$K_0,K_1,K_2$ work for every $\rho$.

For a curve $c$ in the product define the slope
\[
 u(x,t)=\langle S(x,t),U(c(x,t))\rangle_{g(t)+da^2}.
\]
Then $|u|\leq1$ and
\begin{equation}\label{eq:ramps-slope}
 \partial_tu=\partial_s^2u+(k^2+\operatorname{Ric}(S,S))u.
\end{equation}
The lower estimate is deliberately guarded:
\[
 0\leq u\quad\Longrightarrow\quad
 \partial_s^2u-K_2u\leq\partial_tu.
\]
To derive the equation, use the spatial parallelism of $U$, $Q(U,\cdot)=0$,
and $D_tS=D_sH+(k^2+Q(S,S))S$.
Then $u_s=\langle H,U\rangle$ and
$u_{ss}=\langle D_sH,U\rangle$.
The sign guard is necessary because multiplying the Ricci lower bound by
a negative slope reverses its direction.  Positive ramps are preserved by
applying the maximum principle to the exact equation
(\ref{eq:ramps-slope}) first and using the guarded estimate only after the
sign is known.

\section{Local solutions with fixed labels}
\label{sec:ramps-local-solutions}

\begin{theorem}[Local C2 flow, uniqueness, and continuation]
\label{thm:ramps-local}
\uses{def:ramps-shrinking-curve}
\lean{PoincareMT.M63LocalCurveTheory}
\source{MT2007}
\dcref{Claim 19.1, p. 437}
\dcref{Lemma 19.14, pp. 446--447}
\dcref{Lemma 19.24, pp. 454--455}
Let $\gamma:\mathbb R\to M$ be periodic of period $\mathsf p$, $C^2$, and
immersed.  There is $T$ with $a<T<b$ and a $C^2$ shrinking curve $c$ on
$[a,T]$ satisfying $c(x,a)=\gamma(x)$.  For smooth initial labels there is
a solution smooth also in those labels.  Two $C^2$ solutions with the same initial labels agree
on their common closed or half-open interval.  The solution is intrinsically
regular for every positive time.

If a solution on $[a,T)$ has $T\leq b$ and $k\leq K$ there for some finite
$K$, then it extends to $[a,T']$ for some $T'>T$ when $T<b$, and to the
closed endpoint $T'=b$ when $T=b$.  The extension uses the same ambient flow,
not an extension of $F$ beyond $b$.  For a compact parameter space, initial
values and their first two angular jets varying continuously give a family
of solutions on one common short interval with continuous value and jets.
\end{theorem}
\begin{proof}
The local Euclidean model in
\cite[Claim 19.1, p.~437]{MT2007} alone does not settle the issue of
fixed $C^2$ labels in a moving metric.

First embed $M$ smoothly in a Euclidean space by $e$ and choose a smooth
retraction $r:U\to M$ from a neighborhood of $e(M)$, with $r\circ e=1$.
For $z\in U$ and $w$ with $dr_z w\ne0$, define
\[
 A(t,z,w)=|dr_z w|_{g(t)}^{-2},\qquad
 B(t,z,w)=-A(t,z,w)(D(t)^2e)_{r(z)}(dr_z w,dr_z w).
\]
Here $D(t)^2e$ is the coordinate Hessian of the embedding.  These are smooth
on the indicated set and $A>0$.  The ambient gauge equation is
\begin{equation}\label{eq:ramps-gauge}
 q_t=A(t,q,q_x)q_{xx}+B(t,q,q_x).
\end{equation}
On a compact neighborhood of the initial first jet its principal coefficient
has a positive lower bound.  Smooth cutoffs extend the coefficients while
retaining that bound; the short-time solution stays in the neighborhood,
so the cutoffs do not change its equation.  In a constant-speed initial
coordinate the principal coefficient is initially constant.  The periodic
heat equation then supplies the linear inverse: its Fourier modes have
multipliers $e^{-\lambda_j t}$, and an inhomogeneity is integrated against
$e^{-\lambda_j(t-\tau)}$.  Freeze this constant principal part and iterate
the remaining terms on a sufficiently short interval.  Bounds on the
coefficient derivatives and smallness of $A-A(a,q_0,q_{0,x})$ give the
contraction; the same bounds apply to nearby smooth initial curves.

More explicitly, in the unit-speed coordinate of period $L$ the frequencies
are $\omega_j=2\pi j/L$ and $\lambda_j=\omega_j^2$.
The trace state is square-summable with decoder
$a_j\mapsto a_j/(1+\lambda_j)$; this is the periodic $H^2$ norm.
Its first derivative decoder has multiplier
$i\omega_j/(1+\lambda_j)$.  Cauchy--Schwarz bounds both uniformly because
$\sum_j(1+\lambda_j)^{-1}<\infty$.
The forcing space is $L^2([0,T];H^1)$, realized by the decoder
$(1+\lambda_j)^{-1/2}$.  The modewise heat energy inequality controls its
response in $C([0,T];H^2)$ and its high derivative in $L^2([0,T];H^1)$.
Write the physical forcing equation as
\[
 f=(A(t,q,q_x)-1)q_{xx}+B(t,q,q_x),\qquad
 q=e^{t\partial_x^2}q_0+\int_0^te^{(t-r)\partial_x^2}f(r)\,dr.
\]
Multiplication in periodic $H^1$ is bounded: the product rule and
$H^1\subset C^0$ give
$\|fg\|_{H^1}\leq C\|f\|_{H^1}\|g\|_{H^1}$.
On a forcing ball of radius $r$, the difference of the right-hand sides
is bounded by $\kappa_T\|f-\widetilde f\|_{L^2H^1}$, with a bound of the form
\[
 \kappa_T=2\epsilon_T+8K_Pr
       +2K_P\|q_{\mathrm{heat}}\|_{L^2H^3([0,T])}
       +2K_B\sqrt T.
\]
Here $K_P,K_B$ bound the principal and lower coefficient differences in
the chosen neighborhood; $\epsilon_T\to0$ because $A=1$ initially.
The heat energy is integrable for $H^2$ initial data and its integral over
$[0,T]$ tends to zero.  Choose $r$ and then $T$ to make $\kappa_T<1$ and
the zero-forcing residual less than $(1-\kappa_T)r/2$.
The forcing map preserves the ball and contracts, and its fixed point
lies in its interior.  Invertibility of $1-D_f\mathcal R$ by its Neumann
series gives a smooth local branch in the initial state.  Spatial
translations commute with the equation and its locally unique branch;
differentiating that identity gives spatial smoothness for smooth initial
data on the same interval.  This provides the common smooth approximants
used below, not just separate existence times for each approximant.

The ambient solution remains in $e(M)$.  Indeed, for $R=e\circ r$, the
identity $R\circ e=e$ differentiated twice gives
\[
 D^2R(de\,V,de\,V)+DR(D^2e(V,V))=D^2e(V,V).
\]
To see why it keeps the solution in the manifold, define the defect
\[
 \mathcal D(t,z,w)=A(t,z,w)D^2R_z(w,w)
                         +B(t,z,w)-DR_z B(t,z,w).
\]
The displayed Hessian identity gives
$\mathcal D(t,R(z),DR_z w)=0$ at every projected jet.
All actual and projected jets lie in a compact subset of the smooth
coefficient domain.  Thus $\mathcal D$ is Lipschitz there.
For $Z=q-R(q)$, direct differentiation gives
\[
 Z_t-A Z_{xx}=\mathcal D(t,q,q_x),\qquad
 |Z_t-A Z_{xx}|\le C(|Z|+|Z_x|).
\]
If $A\ge\alpha>0$, the product rule and Young's inequality imply
\[
 (\partial_t-A\partial_x^2)|Z|^2
 \le (2C+C^2/\alpha)|Z|^2.
\]
The periodic maximum principle and $Z(\cdot,a)=0$ give $Z=0$.
Put $d=r\circ q$.
The coordinate Hessian chain rule in (\ref{eq:ramps-gauge}) gives
\[
 d_t=|d_x|^{-2}D_xd_x=H+\lambda d_x,\qquad
 \lambda=\frac{\partial_x|d_x|}{|d_x|^3}.
\]
Solve $\partial_t\chi(x,t)=-\lambda(\chi(x,t),t)$ with $\chi(x,a)=x$.
Periodicity and uniqueness give $\chi(x+\mathsf p,t)=\chi(x,t)+\mathsf p$;
its spatial derivative is the positive exponential of the integral of
$-\lambda_x$.  Thus $d(\chi(x,t),t)$ has velocity exactly $H$.

For $C^2$ initial data the limiting step must also recover second spatial
derivatives in the original labels.  Define
\[
 \sigma(x)=\int_0^x|\gamma_y|_{g(a)}\,dy,\quad
 L=\sigma(\mathsf p),\quad \kappa=L/\mathsf p,\quad
 \phi(x)=\sigma(x)/\kappa.
\]
The map $\sigma$ is an increasing $C^2$ diffeomorphism with
$\sigma(x+\mathsf p)=\sigma(x)+L$; $\gamma\circ\sigma^{-1}$ has unit
speed and period $L$.  Mollify its embedded representative, retract to $M$,
and normalize the smooth approximants to constant speed.  Their initial
speeds converge to $\kappa$ after rescaling to period $\mathsf p$.
The preceding heat construction gives them one common interval.
The curvature-square inequality and the first differentiated estimate give
on a smaller common interval
\[
 k_j^2\leq R_0,\qquad |D_sH_j|\leq J/\sqrt{t-a}.
\]
We spell out convergence through $t=a$, since the displayed first-jet
bound alone is singular there.

First retain the gauge curves $q_j$ and their changes of label $\chi_j$
from the common spectral branch.  Convergence of their initial states
gives convergence in $C_tH^2_x$ on the common closed slab.
The value and first derivative therefore converge uniformly, while
the second derivative converges in $C_tL^2_x$.
Their first jets lie in one compact nondegenerate domain; hence the
gauge speeds have uniform positive upper and lower bounds.
The normal speed equation and the common curvature cap give the
same bounds for normal speeds.  Since
\[
 v_j^{\rm normal}(x,t)
   =v_j^{\rm gauge}(\chi_j(x,t),t)\,\partial_x\chi_j(x,t),
\]
there are $0<\ell<U$ such that
$\ell\le\partial_x\chi_j\le U$ throughout the common slab.

The label equation can also be written
$\partial_t\chi_j=\tfrac12\partial_xA_j(\chi_j,t)$.
The chain rule for $A(t,q_j,q_{j,x})$, the compact first-jet bounds
and the $L^2$ bound for $q_{j,xx}$ give
$\|\tfrac12\partial_xA_j\|_{L^2}\le C$ at every time.
Changing variables through the positive label gives
\[
 \|\chi_j(\cdot,t)-\chi_j(\cdot,s)\|_{L^2}
                    \le \ell^{-1/2}C|t-s|.
\]
The displacements $d_j(x,t)=\chi_j(x,t)-x$ are periodic and
uniformly Lipschitz in $x$.  To obtain a uniform time modulus,
let $P_\tau$ be the periodic Gaussian heat operator.  For a
periodic $K$-Lipschitz function $f$,
\[
 \|f-P_\tau f\|_\infty\le C K\sqrt{\tau},
 \qquad
 \|P_\tau f\|_\infty\le C_\tau\|f\|_2.
\]
The first estimate follows by integrating
$|f(x-y)-f(x)|\le K|y|$ against the periodized Gaussian;
the second is Cauchy--Schwarz for its kernel.
Apply them to $d_j(t)-d_j(s)$, choosing $\tau$ first and
$|t-s|$ second.  This proves uniform time equicontinuity.
The zero initial displacement also bounds the displacements.
Arzela--Ascoli gives one uniformly convergent subsequence of
labels on the closed time-period rectangle, retaining
the increment bounds $\ell(y-x)\le\chi(y,t)-\chi(x,t)\le U(y-x)$.
No derivative of this limiting label is used.
Uniform convergence of gauge positions and unit tangents,
composed with these labels, gives the corresponding normal
field limits.

At each positive time, the embedded gauge curvature is an
affine function of $q_{j,xx}$ whose coefficients depend smoothly
on $(t,q_j,q_{j,x})$.  It is consequently Cauchy in $L^2$.
Changing labels retains this property: for any continuous
periodic field $f$,
$\|f\circ\chi_j\|_2\le\ell^{-1/2}\|f\|_2$.
In comparing two compositions, insert one fixed continuous
member $f_N$.  The two outer differences are controlled in
$L^2$, and uniform continuity of $f_N$ controls the middle
difference caused by the labels.  The normal curvature fields
also have a common spatial Lipschitz bound at that positive
time, from $|D_sH_j|\le J/\sqrt{t-a}$ and the speed bounds.
The same Gaussian smoothing estimate therefore upgrades
$L^2$ Cauchy convergence to uniform convergence.

It remains to control the initial trace uniformly.
Write $\mathcal H_j=de(H_j)$, let $v_{j,0}$ be its constant
initial speed, and put $\nu_j=v_{j,0}^{-2}$.
For $\tau=t-a$, the speed equation and its spatial derivative give
\[
 |v_j-v_{j,0}|\le C\tau,\qquad
 |\eta_j|\le C\tau,\qquad
 |\partial_x\eta_j|\le C(\tau+\sqrt{\tau}),
 \quad \eta_j=v_j^{-2}-\nu_j.
\]
Indeed, differentiating $k_j^2+Q(S_j,S_j)$ gives an
upper bound $C(1+\tau^{-1/2})$; integrating it in the
exponential speed formula yields
$|(v_j)_x|\le C(\tau+\sqrt{\tau})$.  Differentiating $v^{-2}$
gives the last estimate.
The actual curvature-vector evolution, pushed through $e$,
now has the divergence form
\[
 \partial_t\mathcal H_j-\nu_j\partial_x^2\mathcal H_j
       =\partial_x F_j+G_j,\qquad
 F_j=\eta_j\partial_x\mathcal H_j,
\]
with
\[
 \|F_j(t)\|_\infty\le C\sqrt{\tau},\qquad
 \|G_j(t)\|_\infty\le C(1+\tau^{-1/2}).
\]
For these estimates the embedding derivatives are bounded on
the compact manifold.  The conversion of the covariant
Laplacian to the fixed parameter contributes
$-v_xv^{-3}\partial_x\mathcal H_j$ and bounded embedding
terms.  The remaining curvature evolution contains at most
one $D_sH_j$ multiplied by bounded curvature or ambient tensors.
The term $\eta_{j,x}\partial_x\mathcal H_j$ introduced by
divergence form is bounded using the estimates above.
Thus no second curvature derivative is required in $G_j$.

Restart Duhamel's formula at a positive time $\sigma$.
The Gaussian derivative has $L^1$ norm at most
$C(\nu_j(t-r))^{-1/2}$.  Since $\nu_j$ has fixed positive
upper and lower bounds, the divergence forcing contributes
at most
$C\int_0^\tau\sqrt r/\sqrt{\tau-r}\,dr\le C\tau$;
the other forcing contributes at most $C(\tau+2\sqrt\tau)$.
For each smooth approximant, let $\sigma\downarrow a$
using its own continuous initial curvature.  The constants
were fixed before $j$, so
\[
 \|\mathcal H_j(t)-P_{\nu_j(t-a)}\mathcal H_j(a)\|_\infty
                        \le C(\tau+\sqrt\tau)
\]
holds uniformly in $j$.  Initial $C^2$ convergence gives
uniform convergence of $\mathcal H_j(a)$.  Their set together
with its limit is compact in $C^0$.  Strong continuity and
the contraction bound of $P_\tau$ give a common initial
modulus on this compact set, by a finite approximation.
The last display supplies that same modulus for
$\mathcal H_j(t)$ at $a$.

On closed slabs after $a$, the first two curvature-jet bounds
and the evolution equation give a common time modulus.
Joining it to the initial modulus proves equicontinuity
on the whole closed slab.  The positive-time uniform limits
and the initial limit therefore assemble to one uniform
limit of embedded curvature through $a$.
The speeds are not selected independently: they obey
\begin{equation}\label{eq:ramps-speed-reconstruction}
 v_j(x,t)=v_j(x,a)\exp\left(-\int_a^t
              (k_j^2+Q(S_j,S_j))(x,\tau)\,d\tau\right).
\end{equation}
They therefore converge uniformly to a strictly positive $v$.
To pass to their spatial derivatives, first fix a time $t>a$.
The second intrinsic curvature-jet bound, the already established
speed-gradient bound and the embedding chain rule bound
$\partial_x^2\mathcal H_j(t)$ uniformly.  Uniform convergence
of $\mathcal H_j(t)$ then implies uniform convergence of its
first spatial derivatives.  Indeed, for a periodic $C^2$ function
$f$ with $\|f''\|_\infty\leq B$, Taylor's formula gives
\[
 \|f'\|_\infty\leq2\|f\|_\infty/h+Bh/2\qquad(h>0).
\]
Apply this to differences, choosing $h$ and then the sequence
indices.  Passing to the limit in the fundamental theorem of
calculus identifies the limiting derivative.
The exact formula
\[
 de(D_sH_j)=v_j^{-1}\partial_x\mathcal H_j
                                  -(D^2e)(S_j,H_j)
\]
now gives convergence of the pushed first intrinsic jets.
Consequently $\partial_x(k_j^2+Q(S_j,S_j))$ converges uniformly
in $x$ at each positive time, by the tensor chain rule.
Its norm is at most $C(1+(t-a)^{-1/2})$, using
$|\partial_xk_j^2|\leq2v_j k_j|D_sH_j|$.
Dominated convergence for these continuous periodic functions
makes their time primitives converge uniformly in both variables.
Differentiating (\ref{eq:ramps-speed-reconstruction}) therefore
gives uniform convergence of $(v_j)_x$ to $v_x$.
In particular $v_x(a)=0$ in the constant-speed coordinate.
Passing to the limit in
\[
 c_x=vS,\qquad D_xc_x=v_xS+v^2H,\qquad c_t=H
\]
identifies the first and second spatial derivatives and the time derivative
of the limiting curve; it is an immersed $C^2$ shrinking curve through the
initial endpoint.  Composing with the fixed $\phi$ restores exactly
$\gamma(x)$ at $a$.

Here is how to identify all higher spatial curvature jets of this
same limit.  On a fixed positive-time closed slab let
$\eta_{i,j}=de(D_s^iH_j)$ for the smooth approximants.  The
smooth-curve derivative estimates bound every fixed order there.
The embedding chain rule, together with the speed and speed-gradient
bounds, therefore bounds $\partial_x^2\eta_{i,j}$ uniformly
in both $x$ and time; it uses only intrinsic orders through $i+2$.
Suppose $\eta_{i,j}$ already converges uniformly on the slab.
Taylor interpolation, applied uniformly in time, makes
$\partial_x\eta_{i,j}$ converge too.  The identity
\[
 \eta_{i+1,j}=v_j^{-1}\partial_x\eta_{i,j}
                           -(D(t)^2e)(S_j,D_s^iH_j)
\]
then gives uniform convergence at order $i+1$.  In the last term
the intrinsic vector is recovered from its embedded image by
$dr$, so its coefficients also converge.  Starting from the
already convergent curvature $\eta_{0,j}$ proves convergence
of every finite order on the same sequence.
Passing to spatial integral identities identifies the limiting
recurrences.  Inductively applying the tangent map of the
retraction identifies their fields with the actual $D_s^iH$.
They are spatially $C^1$ and jointly continuous on the closed
slab.  Joint time differentiability will follow from the fixed-label
smoothing argument below; it is not assumed at this step.

For uniqueness, choose a smooth immersed reference curve sufficiently close
in the first jet to an initial curve.  A nearby curve is a normal graph over
this reference after a small increasing change of parameter.  This is a
labelwise construction: at a self-intersection one uses the nearby parameter
on the reference, rather than selecting a point from its image.  In the
Euclidean embedding the footpoint equation has derivative bounded away
from zero, so the local implicit solutions agree and give a periodic
footpoint map.  If $r_0$ is the reference and $W$ the graph field, the
constraint $\langle W,r_0'\rangle=0$ implies
\[
 \langle W_{xx},r_0'\rangle
       =-2\langle W_x,r_0''\rangle-\langle W,r_0'''\rangle.
\]
Projecting the velocity along the graph tangent cancels the tangential
coefficient and leaves a system $W_t=A W_{xx}+B_0(x,t,W,W_x)$ with
$A>0$.  The constraint also recovers $W_{xx}$ continuously from curvature,
value and first jet at the initial endpoint.  For two such graphs their
difference $Z=W_1-W_2$ satisfies
\[
 |Z_t-A_1Z_{xx}|\leq C(|Z|+|Z_x|).
\]
Indeed the error is $(A_1-A_2)W_{2,xx}+B_1-B_2$;
the coefficients are Lipschitz on the compact graph neighborhood,
and $W_{2,xx}$ is bounded there.  If $A_1\geq\lambda_0>0$,
\[
 (\partial_t-A_1\partial_x^2)|Z|^2
       \leq(2C+C^2/\lambda_0)|Z|^2.
\]
The periodic scalar maximum principle gives equality of the graphs
from their equal initial values.  If their original curves satisfy
$c(x,t)=d(\theta(x,t),t)$, invariance of curvature under the
increasing spatial relabeling and the two pure-normal equations give
$d_x\theta_t=0$.  Immersion implies $\theta_t=0$; its initial
value is $x$.  Thus the original labels agree as well.
Repeating at any common time, and closing the interval of agreement
by continuity, proves uniqueness on closed and half-open intervals.

Continuous dependence of the second jet requires the limiting
construction, in addition to this uniqueness argument.  Here are
the compactness details.  The embedded triples
$(e\gamma,(e\gamma)_x,(e\gamma)_{xx})$ form a compact subset of
the metric space of continuous periodic triples with the uniform norm.
One may work over this image even if the given compact parameter
space is not metrizable.  Near one triple, normalize arclength and
rescale all the nearby curves to the same physical period $L$.
Their constant speeds vary continuously, equal one at the central
triple, and lie in a fixed positive interval.  The normalizing
labels and their first two derivatives vary continuously too,
as follows from the positive-speed primitive and the inverse
derivative formulas.

At successive errors tending to zero, take finite smooth
approximating sets for this compact set of normalized triples.
The union of these sets and the original compact set is compact:
outside any neighborhood of the original set only finitely many
approximating sets remain.  It therefore fits in one strict
spectral branch after the neighborhood is reduced.  The resulting
countable smooth family has one lifespan, one speed interval,
one initial curvature bound and the common initial-trace modulus
proved above.  For any converging sequence of triples from this
family, the five-field construction gives a subsequence converging
in value and both ordinary spatial derivatives on the closed
time slab.  The limit solves the equation with the limiting
initial data.  Uniqueness makes this limit independent of every
subsequence and every choice of smooth approximants.

This proves continuous dependence also at the original triples.
Explicitly, if it failed for triples approaching a fixed triple,
choose from each approximating sequence a smooth member whose
initial triple and solution triple are each within $1/j$ of the
selected member.  The resulting sequence belongs to the same
smooth family, so the preceding subsequence argument contradicts
the asserted failure of convergence.  Restore the original
labels using
\[
 (d\circ\phi)_x=d_x(\phi)\phi',\qquad
 (d\circ\phi)_{xx}=d_{xx}(\phi)(\phi')^2+d_x(\phi)\phi''.
\]
A finite cover of the compact initial image now supplies a common
positive lifespan.  Uniqueness identifies the locally constructed
families on overlaps, giving the required continuous family in
the original labels.

We next justify smooth labels on a positive-time slab $[\tau,s]$.
Normalize the arclength of the slice at $\tau$ by one increasing
$C^2$ map $\psi$, commuting with period translation, and put
$d(x,t)=c(\psi(x),t)$.  The speed of $d$ at $\tau$ is a constant
$v_0>0$.  This normalization does not by itself assert that the
slice is smooth.  Put $q=ed$, $R_0=de(S)$ and
$R_{i+1}=de(D_s^iH)$, with the fields evaluated along $d$.
The preceding spatial-limit argument makes all these fields
continuous through the closed slab.  The chain rule for the embedding gives
\[
 q_x=vR_0,\qquad
 (R_i)_x=v\bigl(B_e(t,q,R_0,R_i)+R_{i+1}\bigr),
\]
where
\[
 B_e(t,z,X,Y)=(D(t)^2e)_{r(z)}(dr_zX,dr_zY).
\]
Also, with
$N_e(t,z,X,Y)=Q_{r(z)}(dr_zX,dr_zX)+|dr_zY|_{g(t)}^2$,
\[
 v(x,t)=v_0\exp\left(-\int_\tau^t
             N_e(u,q(x,u),R_0(x,u),R_1(x,u))\,du\right).
\]
Regard these as functions of $x$ with values in the Banach space
of continuous paths on $[\tau,s]$.  Smooth coefficient substitution
is smooth locally on this space; the time primitive is bounded
linear.  If $q$ and all $R_i$ are $C^m$ as path-valued functions,
the speed formula makes $v$ $C^m$.  The two spatial recurrences
then make $q$ and all $R_i$ $C^{m+1}$.  Starting with continuity,
simultaneous induction proves every spatial order, uniformly
through both endpoints.  In particular each spatial slice of $d$
is smooth, and all its ordinary spatial jets are jointly continuous.

For completeness, this also gives joint smoothness in the open
slab, rather than just smooth slices.  The genuine normal equation
for $q$ is $q_t=\Phi(t,q,q_x,q_{xx})$, with smooth $\Phi$ on the
immersed-jet domain.  In terms of the gauge coefficients above,
\[
 \Phi=Aq_{xx}+B+\tfrac12(\partial_x A)q_x;
\]
here $\partial_x A$ denotes its chain-rule expression in
$(t,q,q_x,q_{xx})$.  On an interior closed subinterval, integrate
this equation in time and differentiate its primitive in $x$.
Continuity of all spatial jets justifies each differentiation and gives
\[
 \partial_t\partial_x^j q
       =\partial_x^j\Phi(t,q,q_x,q_{xx}).
\]
The right-hand side is a smooth expression in finitely many
spatial jets.  Induction on their joint differentiability now
proves joint smoothness of every order.  Composing with $r$
returns to the manifold.  Thus
$c(x,t)=d(\phi(x),t)$ with the single fixed map $\phi=\psi^{-1}$.
The intrinsic jets of the smooth representative, pulled back by
this fixed $C^2$ map, are jointly $C^1$.  Together with their
closed-slab continuity this proves intrinsic regularity of the
constructed solution.  Intrinsic estimates for smooth solutions
transfer through this map.
The estimates used in the preceding limiting construction require
only the smooth-curve dissipation inequalities under an already
supplied curvature cap; their proof is given below.  They do not
assume the later small-turning curvature theorem.

The same regularity holds for every existing $C^2$ solution,
not only the one constructed by approximation.  On a given closed
interval $[a,T_0]$, let $S$ be the supremum of endpoints through
which it is intrinsically regular.  Local existence and uniqueness
give a nonempty set of such endpoints.  Before $S$, use one
constant-speed relabeling at an earlier time $\alpha$; the argument
just given makes that representative smooth on every shorter slab.
The given $C^2$ curve has continuous curvature on the compact
closed cylinder through $S$, hence a uniform curvature cap there.
The smooth derivative bounds give bounds of every order away
from $\alpha$, independent of the shorter upper endpoint, as
well as a speed-gradient bound.  Curvature itself already has
its continuous value at $S$.  The preceding interpolation and
spatial-recurrence argument, now with time tending to $S$ in
place of the approximation index, successively closes every
intrinsic jet at $S$.  It identifies each with the actual spatial
jet of the given terminal slice.

The path-valued recurrence therefore makes that slice smooth
in the fixed normalized labels.  If $S<T_0$, restart its smooth
local solution and use uniqueness to identify it with the given
curve on a larger interval.  Across the join the spatial jets
agree, and the speed primitive splits into the sum of its two
time integrals.  Applying the path-valued recurrence and normal
equation once more proves joint intrinsic regularity across the
join.  This contradicts the definition of $S$, so $S=T_0$.
Restricting to closed subintervals proves the half-open case.
Smooth initial labels retain their smooth representative by
the same argument starting from the initial smooth branch.

Finally suppose $k\leq K$ on $[a,T)$, where $T\leq b$.
Choose $a<\alpha<\tau<T$ and use constant-speed labels at
$\alpha$ as above.  The speed formula gives constants
$0<m\leq v\leq V$ independent of time before $T$.
The smooth-curve derivative estimates, applied on intervals with
upper endpoints increasing to $T$, give bounds for $D_sH$ and
$D_s^2H$ on $[\tau,T)$ independent of those endpoints.
The first-jet estimate near $\alpha$ is integrable in time, so
\[
 I(x,t)=\int_\alpha^t\partial_x(k^2+Q(S,S))(x,u)\,du,
 \qquad v_x=-vI
\]
also has a uniform bound.  Write $r_c=ec$ and $h=de(H)$ in
these labels.  The equations for $r_c,h,v,I$ give uniformly bounded
time derivatives on $[\tau,T)$: the equation for $h$ uses at
most $D_s^2H$, and $I_t=\partial_x(k^2+Q(S,S))$ uses at
most $D_sH$.  They therefore converge in the uniform norm of
continuous periodic functions as $t\uparrow T$.  So does
$g=v_x=-vI$.  The range of $e$ is closed by compactness, and
the limiting speed remains at least $m$.

It remains to recover the tangent and second derivative; bounds
alone would not identify them.  Set $p=(r_c)_x=v\,de(S)$.
Before $T$,
\[
 p_x=g\,de(S)+v^2\bigl(h+(D(t)^2e)(S,S)\bigr).
\]
The right-hand side is bounded.  Applying the Taylor interpolation
estimate used above to $r_c(\cdot,t)-r_c(\cdot,u)$ shows that
$p(\cdot,t)$ is uniformly Cauchy as $t,u\uparrow T$.
Consequently $de(S)=p/v$ converges uniformly.  The displayed
right-hand side now converges uniformly too, since its coefficients
are continuous on a compact set.  Passing to limits in the
spatial integral identities proves that the limiting $r_c$ is
$C^2$, with first derivative $p$ and second derivative the
displayed limit.  The same argument identifies the limiting $g$
as the derivative of the limiting speed.

Retract this terminal curve to $M$.  The unit-length identity
for $S$ passes to the limit, so its velocity has norm $v\geq m$;
the terminal curve is immersed.  The acceleration identity
$D_xc_x=v_xS+v^2H$ identifies $h$ as its actual pushed curvature.
Restoring the fixed original labels gives an immersed $C^2$
endpoint with the required continuous fields.  If $T<b$, restart
the local construction there.  The common limiting curvature
gives the same one-sided normal time derivative for both pieces,
so they join to a $C^2$ shrinking curve.  If $T=b$, the closed
endpoint itself suffices; no equation beyond the ambient time
interval is asserted.
\end{proof}

\section{Positive ramps and uniform estimates}
\label{sec:ramps-positive}

\begin{definition}[Ramp]
\label{def:ramps-ramp}
\lean{PoincareMT.M63IsRampAt}
\uses{def:ramps-shrinking-curve}
For a product $P=M\times C_\rho$, a periodic curve is a ramp at time $t$ if
its slope $u(\cdot,t)$ is strictly positive.  A positive-degree lift is a lift
of the circle coordinate to $\mathbb R$ whose increment over one parameter
period is a positive integer multiple of $\rho$; its degree is that integer.
\end{definition}

The circle component has a $C^2$ lift $\alpha:\mathbb R\to\mathbb R$.
Its periodic increment is $d\rho$ for an integer $d$, and
$\alpha_x=vu$.  Thus positive slope implies $d>0$; it is not a separate
initial hypothesis.  During a continuous evolution the integer $d$ is
constant, since its increment varies continuously.

\begin{proposition}[Global existence for each ramp]
\label{prop:ramps-global}
\lean{PoincareMT.m63RampEstimates_from_predecessors,PoincareMT.M63C2RampExistence}
\uses{def:ramps-ramp,prop:ramps-length-curvature}
For each fixed $\rho>0$, every immersed periodic $C^2$ curve with positive
slope has a unique $C^2$ pure-normal solution on all of $[a,b]$.
Its slope remains positive and its degree is unchanged.  Smooth initial
labels give the smooth-label branch.
\end{proposition}
\begin{proof}
Start with the local solution.  On any compact interval of its existence,
the potential $k^2+Q(S,S)$ in (\ref{eq:ramps-slope}) is bounded.
The maximum principle, with an exponential integrating factor, first gives
$u\geq0$.  Now the guarded lower inequality applies.  If
$m=\min_x u(x,a)>0$, comparison with $m e^{-K_2(t-a)}$ gives
\begin{equation}\label{eq:ramps-slope-lower}
 u(x,t)\geq m e^{-K_2(t-a)}>0.
\end{equation}
For general $C^2$ labels, make these calculations on positive-time smooth
relabelings and pass to $a$ using endpoint continuity.

Set $w_\varepsilon=h_\varepsilon/u$.  The quotient rule gives
\[
 (w_\varepsilon)_t-(w_\varepsilon)_{ss}
       -2\frac{u_s}{u}(w_\varepsilon)_s
 =\frac{h_{\varepsilon,t}-h_{\varepsilon,ss}}u
       -\frac{h_\varepsilon}{u^2}(u_t-u_{ss}).
\]
Since $h_\varepsilon\geq k$, the cubic term becomes
$(k^3-h_\varepsilon k^2)/u\leq0$.  Consequently
\begin{equation}\label{eq:ramps-forced-quotient}
 (w_\varepsilon)_t\leq(w_\varepsilon)_{ss}
   +2\frac{u_s}{u}(w_\varepsilon)_s
   +(C_1+K_2)w_\varepsilon+\frac{C_1}{u}.
\end{equation}
The final forcing term cannot be dropped; compare
\cite[pp.~8--9, corrected ramp comparison]{MT2015Correction}.
If $w_\varepsilon(x,a)\leq R_0$, (\ref{eq:ramps-slope-lower}) and the
maximum principle on any existing interval $[a,T]$ give
\[
 w_\varepsilon(x,t)\leq
 \left(R_0+\frac{C_1e^{K_2(T-a)}}m(t-a)\right)
                e^{(C_1+K_2)(t-a)}.
\]
The comparison function has derivative at least its right-hand linear
reaction plus the constant forcing bound.  Taking $\varepsilon=1$, and
using $0<u\leq1$, bounds $k\leq h_1\leq w_1$ up to any finite terminal
time, with a bound obtained by replacing $T$ by $b$.  Curvature-bounded
continuation closes the maximal interval at $b$ and rules out an earlier
terminal time.  Uniqueness identifies all local extensions.
\end{proof}

The lower bound $m$ may tend to zero with $\rho$.  This argument proves
existence for each circumference; it supplies no circumference-independent
curvature cap.  The next estimate supplies a different uniform bound.

\begin{theorem}[Uniform short-scale derivative estimates]
\label{thm:ramps-uniform-jets}
\lean{PoincareMT.m63RampEstimates_from_predecessors,PoincareMT.M63UniformDerivativeEstimates}
\uses{prop:ramps-global,prop:ramps-length-curvature}
\source{MT2007}
\dcref{Lemma 19.24, p. 455}
\dcref{Lemma 19.58 and proof, pp. 481--496}
Given nonnegative initial bounds $L_0,\Theta_0$, there are
$0<\delta_0<1$, a radius $r_0=1$, and nonnegative constants $C_i$ chosen
before $\rho$, the curve, the reference time, and the scale.  They apply to
every existing $C^2$ shrinking curve on the full interval $[a,b]$ in
$M\times C_\rho$, with $L(a)\leq L_0$, $\Theta(a)\leq\Theta_0$ and
$0<\rho<1$, as follows.  Put $\delta=\delta_0r_0^2$.  If
\[
 a\leq s,\quad L(s)\geq r,\quad 0<r<1,
\]
every parameter subarc of length at most $r$ has total curvature at most
$\delta$, and $s+\delta r^2<b$, then for $s<t<s+\delta r^2$ and every
$i\geq0$,
\[
 |D_s^iH|^2\leq {C_i\over(t-s)^{i+1}}.
\]
The supporting curvature estimate uses $0<r\leq r_0$, subarc total curvature
at most $\delta_0$, and $s<t\leq s+\delta_0r^2\leq b$; it gives the base
curvature bound
\[
 |H|^2\leq {2\over t-s}.
\]
The higher-jet estimate is the separate assertion above.  The reference time
is excluded from both inequalities.  Here a parameter subarc is
$[\alpha,\beta]$ with $\alpha\leq\beta\leq\alpha+\mathsf p$, its length
is $\int_\alpha^\beta v\,dx$, and its total curvature is
$\int_\alpha^\beta kv\,dx$.  The derivatives $D_s^iH$ are full covariant
derivatives, without normal projection at each step.  No slope hypothesis
is imposed.
\end{theorem}
\begin{proof}
We first prove the curvature bound in smooth labels.  Choose $K\geq0$
bounding $\operatorname{Rm}$, $D\operatorname{Rm}$, $Q$, $DQ$ and $D^2Q$
on all the products.  Such a common bound exists because their tensors are
horizontal pullbacks from $M$.  The integral estimates give a common $E$
with $L(t),\Theta(t)\leq E$ for $a\leq t\leq b$, for example
\[
 E=(L_0+\Theta_0)e^{(C_1+K_2)(b-a)}.
\]
The construction of $\delta_0$ depends on both initial bounds through $E$.

Suppose temporarily that $k^2\leq2/(t-s)$ up to some time $T$.
For any fixed-label arc $J$ of initial length at most $r$, the speed
equation gives both $L_J(t)\leq e^{K(t-s)}L_J(s)$ and
\begin{equation}\label{eq:ramps-arc-loss}
 L_J(s)-L_J(t)
 \leq Ke^Kr(t-s)+2\sqrt2 E\sqrt{t-s},\qquad t-s\leq1.
\end{equation}
Indeed $\int_Jk^2\,ds\leq\sqrt{2/(t-s)}\Theta_J(t)$, and integrating
this inverse square-root bound gives the second term.  This loss estimate
remains integrable at the restart time.

Fix a $C^2$ cutoff $0\leq\psi\leq1$ that is one on $[-3/8,3/8]$,
supported in $(-9/20,9/20)$, with $|\psi'|\leq P_1$ and
$|\psi''|\leq P_2$.  For a centered arc $J$ and its fixed center label $x_0$,
write $\sigma(x,t)=\int_{x_0}^xv(y,t)\,dy$ and set
\[
 I(t)=\int_J\psi(\sigma(x,t)/r)k(x,t)\,ds.
\]
Keep both sides of $J$ at least $9r/20$ from the center, so the cutoff
vanishes near its endpoints.  Differentiating its regularized version and
integrating twice by parts gives
\begin{align*}
 I_\varepsilon'\leq{}&
 \left(\frac{P_2}{r^2}+\frac{P_1}{r}
       \left(KL_J+\sqrt{\frac2{t-s}}\Theta_J\right)\right)
                 (\Theta_J+\varepsilon L_J)\\
 &+(C_1+K)I_\varepsilon+C_1L_J,
\end{align*}
where here $C_1=C_1(K,K,K)$.  The term involving $P_1$ comes from
$|\sigma_t|\leq KL_J+\sqrt{2/(t-s)}\Theta_J$.
The final $C_1L_J$ is the moving-metric correction.  With
$L_J\leq L_*$ and $\Theta_J\leq\Theta_*$, integrate first and then send
$\varepsilon$ to zero to obtain
\begin{equation}\label{eq:ramps-local-turning}
 I(t)\leq I(s)+D_0(t-s)+2D_1\sqrt{t-s},
\end{equation}
where
\begin{align*}
 D_0&=(P_2/r^2+(P_1/r)KL_*+C_1+K)\Theta_*+C_1L_*,\\
 D_1&=(P_1/r)\sqrt2\Theta_*^2.
\end{align*}
The regularization error is at most $\varepsilon L_J(t)$, so this holds
at both endpoints.  This corrects the localized estimate in
\cite[Claim 19.61, pp.~484--485]{MT2007}; see
\cite[p.~9]{MT2015Correction}.

Choose a positive $\delta_0<1$ so small that
\begin{align*}
 Ke^K\delta_0+2\sqrt2E\sqrt{\delta_0}&\leq1/20,&
 e^{K\delta_0}&\leq3/2,\\
 2(\delta_0+D\sqrt{\delta_0})&<h,
\end{align*}
where
\begin{align*}
 D&=(P_2+P_1KE+C_1+K)E+C_1E+2P_1\sqrt2E^2,\\
 B_*&=(272+17C_0(K,K,K))(114+10K)
                       +2048+992K+160000K^2,\\
 h&=\min\{1/4,(8\sqrt{B_*}+1)^{-1}\}.
\end{align*}
Under the temporary curvature bound, (\ref{eq:ramps-arc-loss}) loses at
most $r/20$ before $s+\delta_0r^2$.  A later arc of length at most $r/2$
must have had initial length at most $r$: otherwise its initial prefix of
length $r$ would still have length at least $19r/20$, a contradiction.
Bisect the arc in its initial arclength coordinate.  Enlarge each half
to an initial arc of length $r$ centered at that half's initial midpoint.
The outer half-lengths stay at least $r/2-r/20=9r/20$, while the target
half stays within $e^{K\delta_0}r/4\leq3r/8$ of its center.
The cutoff therefore has the required support and is one on the target.
Applying (\ref{eq:ramps-local-turning}) to both halves gives
\[
 \Theta_J(t)\leq2(\delta_0+D\sqrt{\delta_0})
 \quad\hbox{whenever } L_J(t)\leq r/2.
\]
The fractions $3/8$, $9/20$ and the loss $1/20$ ensure that the cutoff
boundary remains inside its controlled arc.

If the curvature bound failed, compactness of a period cylinder would
give a first time $T_*>s$ and a label $x_*$ at which
$(T_*-s)k^2(x_*,T_*)=2$.  Put $\tau=T_*-s$.
The first differentiated curvature estimate under the temporary cap gives
$|D_sH|^2\leq B_*/\tau^2$ on the interior time interval
$(s+3\tau/4,T_*)$.  On that interval $k\leq2/\sqrt\tau$, so integrating
$|\partial_sk^2|\leq2k|D_sH|$ gives a scalar oscillation estimate.
Continuity of curvature and arc length then passes that estimate to
$T_*$, even if this is the final time $b$:
\[
 |k^2(y,T_*)-k^2(x,T_*)|
 \leq\frac{4\sqrt{B_*}}{\tau^{3/2}}L_{[x,y]}(T_*).
\]
The differentiated estimate uses a time-weighted sum of $|D_sH|^2$ and
$k^2$, detailed below.  This passage does not require a first curvature
derivative at the closed final slice.

Starting at $x_*$, take a forward arc of length $h\sqrt\tau$.
It exists within a period because (\ref{eq:ramps-arc-loss}) keeps the
total length at least $19r/20$, whereas $h\sqrt\tau\leq r/4$.
The oscillation bound and $4\sqrt{B_*}h\leq1$ give
$k\geq1/\sqrt\tau$ throughout that arc.  Its total curvature is at least
$h$, contradicting the transported upper bound
$2(\delta_0+D\sqrt{\delta_0})<h$.  This proves the curvature bound,
including its upper endpoint.

For higher derivatives put $q_i=|D_s^iH|^2$.
Commuting $D_t$ and $D_s$, differentiating the curvature equation, and
using metric compatibility produces $-2q_{i+1}$ and lower-order
contractions of curvature jets and ambient curvature derivatives.
The contraction containing the highest derivative is absorbed by
$2ab\leq a^2+b^2$, leaving one negative copy of $q_{i+1}$.
The tensors $D^j\operatorname{Rm}$ and $D^jQ$ have product-independent
bounds of every fixed order.  One can track the differentiations without
assuming a normal connection formula: if $J_i=D_s^iH$,
$q=k^2+Q(S,S)$, $\mathcal C=\partial_tD$ and
$E_i=D_tJ_i-J_{i+2}$, the exact recursion is
\[
 E_{i+1}=D_sE_i+qJ_{i+1}+R(H,S)J_i+\mathcal C(S,J_i),
\]
starting from
$E_0=q_sS+2qH+R(H,S)S+\mathcal C(S,S)$.
Here is the estimate for this recursion.  Set $V_0=S$ and
$V_{r+1}=J_r$.  Differentiation and Pascal's identity give
\[
 \langle E_m,J_m\rangle
 =\sum_{r=0}^{m+1}\binom{m+2}{r+1}
       (\partial_s^r q)\langle V_{m+1-r},J_m\rangle
       +\mathcal A_m.
\]
The remainder $\mathcal A_m$ is the sum of the actual Riemann
and differentiated Ricci contractions from the recursion.
Its initial terms are
$\operatorname{Rm}(H,S,J_m,S)$ and
$-2(D_SQ)(S,J_m)+(D_{J_m}Q)(S,S)$.
At each step differentiate every tensor and every non-test
vector slot, then add the displayed curvature and
connection-variation commutators.  Derivatives of the test
field cancel when recovering the paired vector error.
Give $V_r$ weight $r$ and each tensor derivative weight one.
Every resulting term has non-test weight $m+1$; the tensor
orders are at most $m$ for Riemann and $m+1$ for Ricci.
The sum of absolute coefficients is bounded recursively by
$b_0=4$, $b_{m+1}=(m+4)b_m+4$: there are at most $m+4$
places to differentiate an existing term and four new
commutator contributions.

Suppose $|J_i|\le L$ for $i<m$, with $L\ge1$, and all those
ambient tensors have unit-input norm at most $K$.
The weight bound allows at most one factor $V_{m+1}=J_m$
in each ambient term; if it occurs, all other non-test vector
weights are zero.  Tensor multilinearity therefore gives
\[
 |\mathcal A_m|\le b_m K L^{m+3}(1+|J_m|)|J_m|.
\]
Leibniz differentiation of $q=|H|^2+Q(S,S)$ similarly gives,
with $D_0=(m+3)!\,K L^{m+3}+2^{m+1}L^2$,
\[
 |\partial_s^r q|\le D_0\ (r<m),\qquad
 |\partial_s^m q|\le D_0(1+|J_m|),
\]
\[
 |\partial_s^{m+1}q-2\langle H,J_{m+1}\rangle|
                                  \le D_0(1+|J_m|).
\]
The only remaining highest-jet term multiplies
$\langle S,J_m\rangle$.  Differentiating
$\langle H,S\rangle=0$ exactly $m$ times shows
\[
 \langle J_m,S\rangle
   =-\sum_{i=0}^{m-1}\binom mi\langle J_i,J_{m-1-i}\rangle,
 \qquad |\langle J_m,S\rangle|\le2^m L^2.
\]
It thus has a bound using only lower jets, rather than
$|J_m|$.  Bounding the other rows and summing their binomial
coefficients gives, for $m\ge2$,
\[
 |\langle E_m,J_m\rangle|
       \leq T_m|J_{m+1}|+D_m(|J_m|+|J_m|^2),
\]
where one may take
$T_m=2^{m+1}L^3$ and
$D_m=2^{m+2}D_0L+b_mKL^{m+3}$.
Insert this in
\[
 (\partial_t-\partial_s^2)|J_i|^2
       =-2|J_{i+1}|^2-2Q(J_i,J_i)+2\langle E_i,J_i\rangle.
\]
The inequalities $2T_i|J_{i+1}|\leq|J_{i+1}|^2+T_i^2$ and
$2|J_i|\leq1+|J_i|^2$ give the claimed dissipation with finite
constants; orders zero and one use the corrected evolution and its first
derivative separately.  Induction therefore reduces the estimate,
on a positive-time window, to finitely many inequalities
\[
 (\partial_t-\partial_s^2)q_i
       \leq-q_{i+1}+Aq_i+B/(t-\alpha)^i,\qquad 0\leq i\leq N,
\]
with $q_0\leq R$ and time-window length at most $H_*$.
The constants depend on lower jets, the derivative order and ambient
bounds, but not on circumference.  To retain this fact at small
time scales, the same argument is performed with normalized
jets $\lambda^{i+1}|J_i|$, where $0<\lambda\le1$.
Weight counting gives the precise form
\[
 \lambda^{2i+4}(\partial_t-\partial_s^2)q_i
 \le-\lambda^{2i+4}q_{i+1}
                      +C_i\lambda^{2i+2}q_i+C_i.
\]
Every additional nonnegative power of $\lambda$ from an
ambient term is at most one.  For $i\ge2$ one can take
$C_i=2K+4D_i+T_i^2$ using the normalized lower-jet bound $L$.
Orders zero and one are kept separate.  In particular, under
$q_0\le R$, direct differentiation of the corrected equation gives
\[
 (\partial_t-\partial_s^2)q_1
 \le-q_2+(14R+10K+1)q_1+
       4R^3+2RK(7R+6)+\{K(46R+32)\}^2.
\]
More explicitly, before Young's inequality the upper bound is
\[
 -2q_2+(14k^2+10K)q_1
 +K(20k^2+26k+6)|J_1|+4k^3|J_2|
 +2k^2K(1+5k+2k^2).
\]
Here $4k^3|J_2|\leq q_2+4k^6$, whereas the linear
$|J_1|$ term is at most
$q_1+\{K(46R+32)\}^2$, since $k\leq R+1$.
This explains both the remaining dissipation and the last square.

For completeness, write $\mathcal L=\partial_t-\partial_s^2$,
$A_1=14R+10K+1$ and let $G_1$ be the constant forcing in
the first-jet inequality.  The order-zero equation and
$|D_sH|^2=|(D_sH)^\perp|^2+k^4$ give
\[
 \mathcal Lq_0\leq-2q_1+D_0',\qquad
 D_0'=4R^2+C_0(K,K,K)(2R+1).
\]
On a window of length at most $H_*$ put $\Lambda=1+A_1H_*$.
For a positive restart $\sigma$ the function
$W_1=(t-\sigma)q_1+\Lambda q_0$ satisfies
\[
 \mathcal LW_1\leq H_*G_1+\Lambda D_0',\qquad
 W_1(\sigma)\leq\Lambda R,
\]
because the coefficient of $q_1$ is
$1+A_1(t-\sigma)-2\Lambda\leq0$.
Comparison followed by $\sigma\downarrow\alpha$ therefore gives
\[
 q_1(t)\leq\frac{\Lambda R}{t-\alpha}
                         +H_*G_1+\Lambda D_0'.
\]
In the preceding oscillation argument use
$\alpha=s+\tau/4$, $R=8/\tau$ and $H_*=\tau\leq1$.
Then $\Lambda\leq114+10K$,
$\tau^2D_0'\leq256+17C_0(K,K,K)$ and
$\tau^3G_1\leq2048+992K+160000K^2$.
For $t>s+3\tau/4$ we have $t-\alpha\geq\tau/2$;
substitution gives exactly $q_1(t)\leq B_*/\tau^2$.
This derivation is independent of the higher-order induction.

The finite comparison has an explicit cancellation.  Set
$C=N+1+AH_*$ and, for a positive restart $\sigma>\alpha$, form
\[
 W=\sum_{i=0}^N C^{N-i}(t-\sigma)^i q_i.
\]
The differentiated weights and the $Aq_i$ terms are absorbed by the
negative $q_i$ term from the preceding index, since
$i+A(t-\sigma)\leq C$.  Also $(t-\sigma)^i/(t-\alpha)^i\leq1$.
Thus
\[
 (\partial_t-\partial_s^2)W\leq
 D_*:=AC^NR+B\sum_{i=0}^NC^{N-i},\qquad W(\sigma)\leq C^NR.
\]
Comparison and then $\sigma\downarrow\alpha$ give
$q_N(t)\leq(C^NR+D_*(t-\alpha))/(t-\alpha)^N$.
No higher initial trace is used.  For a desired time $t>s$, put
$\lambda=\sqrt{t-s}$ and $\alpha=s+(t-s)/4$.
Work on an interior window from $\alpha$ to a point beyond $t$.
The cap gives $q_0\le8/\lambda^2$ there, and induction gives
one bound for all $\lambda^{i+1}|J_i|$, $i<N$.
The normalized dissipation just proved yields the finite
comparison with $A=\mathcal A/\lambda^2$,
$B=\mathcal A/\lambda^4$, $R=8/\lambda^2$ and
$H_*=\lambda^2$, after enlarging a single constant
$\mathcal A$ to cover all orders through $N$.
Indeed its constant forcing
$C_i/\lambda^{2i+4}$ is at most
$\mathcal A/(\lambda^4(u-\alpha)^i)$ when
$0<u-\alpha\le\lambda^2$.
The weight $N+1+AH_*$ is independent of $\lambda$, while
$t-\alpha=3\lambda^2/4\ge\lambda^2/2$.
The finite comparison therefore gives
$q_N(t)\le C_N/\lambda^{2N+2}=C_N/(t-s)^{N+1}$.
This is the time-weighted mechanism of
\cite[Theorem 3.1, pp.~4--6]{AltschulerIMA822}, with the moving ambient
contractions retained.

Finally use the fixed positive-time relabeling from
Theorem~\ref{thm:ramps-local}.  Length, turning and the full $D_s^iH$ are
invariant under orientation-preserving relabeling.  Restrictions to interior
slabs and endpoint continuity transfer the curvature estimate to $C^2$
labels; the higher estimates are needed only at interior times.
All constants precede the product, curve and time-scale choices.
\end{proof}

\section{Flattened geodesic polygons and canonical graph ramps}
\label{sec:ramps-polygons}

Let $N>0$.  Divide the parameter circle into cells of length
$\ell_N=\mathsf p/N$.  A geodesic $N$-polygon at time $t$ consists of cyclic
vertices and minimizing geodesic sides, with the same metric and connection
$g(t),D(t)$ used to choose every side.  Define
\[
 f_N(x)=\begin{cases}
 \exp\bigl(-1/(1-\cos(Nx))\bigr),&1-\cos(Nx)>0,\\
 0,&1-\cos(Nx)=0,
 \end{cases}\qquad
 \psi_N(x)=\frac{\ell_N f_N(x)}{\int_0^{\ell_N}f_N(y)\,dy}.
\]
The profile is periodic with period $\ell_N$, nonnegative, positive in the open
cell, symmetric about its midpoint, flat to every order at cell endpoints,
and normalized by
\[
 \int_0^{\ell_N}\psi_N(x)\,dx=\ell_N.
\]
Its primitive $\Phi_N$ is strictly increasing and satisfies
$\Phi_N(x+\ell_N)=\Phi_N(x)+\ell_N$ and
$\Phi_N(x+\mathsf p)=\Phi_N(x)+\mathsf p$.  Thus it descends to a
homeomorphism of the parameter circle; flat vertices mean that we claim a
homeomorphism, not a diffeomorphism.

Flatten each geodesic side by composing it with $\Phi_N$ in its cell.  The
result is a smooth periodic curve.  On a side with speed $q_j$, its velocity
is exactly
\[
 \psi_N(x)\,\dot\sigma_j(\Phi_N(x)-\text{cell origin}),
\]
so the horizontal speed is side speed times the profile, not side length
times the profile.  Put the flattened curve into $M\times C_\rho$ using the
circle coordinate whose velocity is $(\rho/\mathsf p)U$.  The resulting
canonical graph ramp has degree one, and
\begin{align}
 L_{P}(\text{ramp})&\leq L_M(\text{flattened polygon})+\rho,\\
 \Theta_{P}(\text{ramp})&\leq N\pi.
\end{align}
Constant sides are allowed.  The graph is immersed even at a stopped base
vertex, since its vertical speed is positive.

\begin{proof}[Profile and graph estimates]
The function $e^{-1/y}$ extended by zero for $y\leq0$ is smooth and flat
at zero.  Composition with $1-\cos(Nx)$ proves flatness at every cell
endpoint.  Positivity in the open cell makes the normalizing denominator
positive and the primitive strictly increasing on every nontrivial
interval.  Periodicity and its normalized integral give all the translation
identities for $\Phi_N$.  It maps each cell onto itself and has a continuous
inverse on the circle.  Every positive-order derivative of a composed side
vanishes at a vertex: each chain-rule term contains a positive-order
derivative of $\Phi_N$.  The sides therefore glue smoothly, even when their
directions disagree.

On a nonconstant side choose its parallel unit tangent $E$, lifted
horizontally to the product, and put $B=\rho/\mathsf p$.
The graph has velocity $q_j\psi_NE+BU$ and speed
$v=\sqrt{q_j^2\psi_N^2+B^2}$.  Hence
\[
 v\leq q_j\psi_N+B,\qquad
 u=B/v>0,\qquad
 kv=\frac{Bq_j|\psi_N'|}{B^2+q_j^2\psi_N^2}.
\]
The last identity follows by differentiating the unit vector
$(q_j\psi_NE+BU)/v$ in the parallel frame $(E,U)$.
Its turning angle from $U$ is $\arctan(q_j\psi_N/B)$.
On one cell the profile increases to its midpoint maximum and decreases
back to zero, so
\[
 \int_{\text{cell}}k\,ds
 =2\arctan\left(\frac{q_j\max\psi_N}{B}\right)\leq\pi.
\]
A constant side has vertical tangent and contributes zero; it does not
require choosing $E$.  Summing proves $\Theta\leq N\pi$.
Integrating the speed bound and the normalization of $\psi_N$ gives
$L_P\leq\sum_jq_j\ell_N+\rho$; the sum is precisely the flattened
polygon's length.  The angular coordinate increases by $\rho$ over one
period, proving degree one.  This is the construction of
\cite[Claims 19.19--19.22, pp.~450--453]{MT2007}.
\end{proof}

\section{Approximation of raw null families}
\label{sec:ramps-family-approximation}

For the rest of the chapter let $\dim M=3$.
Let $\Gamma:S^2\to\Lambda^1M$ be a continuous family of $C^1$ free loops,
each pointwise nullhomotopic.  Here $\Lambda^1M$ is the actual fixed-label
$C^1$ loop space, with topology given by uniform convergence of values and
angular first jets.  For every tolerance $\zeta>0$, there are an integer
$N>0$, a family of geodesic $N$-polygons, and a continuous family of smooth loops
$\widetilde\Gamma$ with the following properties:
\begin{enumerate}
\item each polygon samples $\Gamma(z)$ at the same cell points;
\item $\widetilde\Gamma$ is pointwise nullhomotopic and is homotopic to
      $\Gamma$ through the displayed loop space;
\item the flattened polygon is exactly the periodic representative of
      $\widetilde\Gamma$, and its first and second angular jets are
      continuous in $(z,x)$;
\item for every $z$,
\[
 0\leq L(\Gamma(z))-L(\widetilde\Gamma(z))<\zeta,\qquad
 |A(\widetilde\Gamma(z))-A(\Gamma(z))|<\zeta,
\]
where $A(\gamma)=A_{g(a)}(\gamma)$ is the infimum of
areas of the admissible spanning disks from the filling-area chapter.
All lengths in this approximation statement use $g(a)$.
\end{enumerate}
The same $N$, polygon count, homotopy, and family are retained in every
later construction.  The approximation is not reselected separately for
each circumference or each family member.

\begin{proof}[Construction of the approximation]
Compactness of $S^2$ and continuity into the $C^1$ loop space give a uniform
bound on angular speed and a common modulus of continuity of the first
jets in finitely many ambient charts.  Choose equal cells sufficiently
small that each original segment lies in a common geodesically convex
neighborhood.  Join its endpoints by the unique minimizing geodesic.
Minimality gives a nonnegative length loss on each cell.
Here is a quantitative reason that the sum of the losses tends
uniformly to zero.  Let $S$ bound all original speeds and put
$L_*=S+1$.  Fix $0<\sigma\leq1$ with
$7\sigma L_*\mathsf p<\zeta$, and set $K_*=1+\sigma$,
$d=\sigma L_*$.  In sufficiently small charts, a linear change
of coordinates freezes the metric at the chart center and gives
\[
 |\gamma_x|_g\leq K_*|b(x)|,\qquad
 |q(t)-q(s)|\leq K_*\operatorname{dist}_g(\gamma(t),\gamma(s)),
\]
where $q$ is the coordinate curve and $b=q'$.
Continuity of the first jets makes $|b(u)-b(s)|\leq d$
on small cells.  A Lebesgue number for the finite cover of the
compact parameter-circle product makes all three estimates
valid on every sufficiently short cell, in an appropriate chart.
For a cell of parameter length $\ell$, write $I$ for its
curve length and $D$ for the distance between its endpoints.
Integrating $b-b(s)$ and the speed estimate gives
\[
 \ell|b(s)|\leq K_*D+d\ell,\qquad
 I\leq K_*(|b(s)|+d)\ell.
\]
Since $D\leq I\leq S\ell$,
\[
 0\leq I-D\leq(K_*^2-1)D+2K_*d\ell
                         \leq7\sigma L_*\ell.
\]
Sum over the $N$ cells to obtain total loss less than $\zeta$.

The geodesic interpolator is smooth in its two endpoints.  Compose its
cell parameter with $\Phi_N$ to obtain $\widetilde\Gamma$.
Its length equals the sum of side lengths by the preceding calculation.
Both the raw point and its flattened replacement stay within a short
cell neighborhood, so their distance tends uniformly to zero as $N\to\infty$.
The first two angular jets of the replacement are, within a cell,
\[
 \widetilde\Gamma_x=\psi_N\dot\sigma_j,\qquad
 D_x\widetilde\Gamma_x=\psi_N'\dot\sigma_j,
\]
because the side is a geodesic.  Both jets vanish at vertices and depend
continuously on the endpoints.  They glue continuously in $(z,x)$.
Smoothness in $z$ is neither assumed nor asserted.

For the family homotopy, use a controlled smooth local contraction
$C(\tau,p,q)$, defined for nearby pairs, with
$C(0,p,q)=q$ and $C(1,p,q)=p$.  Uniform closeness puts
\[
 H(\tau,z)(x)=C(\tau,\widetilde\Gamma(z)(x),\Gamma(z)(x))
\]
in its domain for every $(\tau,z,x)$.  Its angular derivative is
$d_pC\,\widetilde\Gamma_x+d_qC\,\Gamma_x$, which is jointly continuous.
Thus $H$ is a homotopy in the $C^1$ loop space, with exact endpoint loops,
and transfers pointwise nullhomotopy.  No continuous family of spanning
disks is required.

The filling-area comparison uses
Proposition~\ref{prop:filling-uniform-collars}: for every total
boundary-length bound $\Lambda$ and $\eta>0$
there is $d>0$ such that two pointwise $d$-close null $C^1$ loops of combined
length at most $\Lambda$ admit attachment of a parametrized collar of area
less than $\eta$ to any disk on either boundary.  Choose
$\Lambda=2\max_zL(\Gamma(z))$ and $\eta=\zeta/2$, and refine the same
sampling mesh to meet this closeness threshold.  For a fixed $z$, take a
disk on $\Gamma(z)$ with area less than $A(\Gamma(z))+\zeta/2$ and attach
the collar.  It is an admissible disk on $\widetilde\Gamma(z)$, so
$A(\widetilde\Gamma(z))<A(\Gamma(z))+\zeta$.
Interchanging the loops gives the reverse strict inequality.  The disks
may be selected separately for each $z$; their role is the numerical area
bound, while the displayed $H$ supplies continuity of the family.
This proves all the claims using one common $N$.
Compare \cite[Lemma 19.17, pp.~449--453]{MT2007}.
\end{proof}

The family length supremum is finite and attained because $S^2$ is compact.
Choose one initial bound $B$ satisfying
\[
 B>\max\{\sup_zL(\Gamma(z))+1,\,N\pi\}.
\]
The canonical ramps then have initial length at most
$\sup_zL(\Gamma(z))+1$ whenever $\rho<1$, and initial total curvature at
most $N\pi$.  Uniform derivative constants are chosen from $B$ before
$\rho$ and before the member $z$.

\section{The solution family and cross-chapter contract}
\label{sec:ramps-family-solutions}

For each fixed $\rho>0$, evolve every canonical ramp in the same product flow.
The result is a continuous family $c_z(x,t)$ of actual shrinking curves with
\begin{align*}
 c_z(x,a)&=\operatorname{ramp}_\rho(\widetilde\Gamma(z))(x),\\
 u_z(x,t)&>0,\qquad \deg(c_z(\cdot,t))=1,\\
 \pi_Mc_z(\cdot,t)&\in\Lambda^1M,\qquad
 \pi_Mc_z(\cdot,a)=\widetilde\Gamma(z).
\end{align*}
The projected loops remain nullhomotopic by the actual continuous homotopy
in the solution-family record.  Values and angular velocities are continuous
in $(z,x,t)$; the initial first and second jets are those of the fixed
approximation.  Length and total curvature obey, uniformly for $0<\rho<1$,
\[
 L_z(t)\leq (\sup_zL(\Gamma(z))+1)e^{K_2(t-a)},
\]
and
\[
 \Theta_z(t)+L_z(t)\leq
 (\sup_zL(\Gamma(z))+1+N\pi)
 e^{(C_1+K_2)(t-a)}.
\]

\begin{proof}[Continuity and retention of the initial family]
Fix $\rho>0$ and choose the unique smooth-label ramp solution for each
$z$.  The initial first and second angular jets are continuous on the
compact space $S^2\times\mathbb R/(\mathsf p\mathbb Z)$.
Their speeds have a positive common minimum, because the vertical speed
is $\rho/\mathsf p$.  The initial slopes have a positive common minimum
$m_\rho$, and $h_1/u$ has a finite common initial maximum.
Proposition~\ref{prop:ramps-global} therefore supplies one curvature cap
for all $z$ on $[a,b]$, for this fixed $\rho$.

Continuous dependence from Theorem~\ref{thm:ramps-local} first gives a
common continuous family on a short interval.  To continue continuity
through the final endpoint, embed the product by $e$ and write
$p=(ec)_x$, $\widehat H=de(H)$ and $g_1=v_x$.
For every fixed $\tau>a$, the common curvature cap and the first two
intrinsic derivative bounds give common time Lipschitz constants for
$ec,p,\widehat H,v,g_1$ on $[\tau,b]$.
For example, with $q=k^2+Q(S,S)$,
\[
 v_t=-qv,\qquad (g_1)_t=-qg_1-vq_x,\qquad
 q_x=v\{2\langle H,D_sH\rangle+(D_SQ)(S,S)+2Q(H,S)\}.
\]
The speed gradient is bounded by integrating this linear equation;
the possible $1/\sqrt{t-a}$ first-jet bound is integrable.
The equations for $p$ and $\widehat H$ involve respectively
$D_sH$ and $D_s^2H$ with bounded lower terms.  They give the
asserted bounds after $\tau$, with endpoint values obtained by
continuity.  The ordinary second derivative is recovered by
\[
 (ec)_{xx}=(D^2e)(c_x,c_x)+v^2\widehat H+(g_1/v)p.
\]
The right-hand side is a continuous function on one compact
coefficient domain, with $v$ bounded below.  Uniform continuity
on that domain supplies a common time modulus for the triple
$\mathcal J_z=(ec_z,(ec_z)_x,(ec_z)_{xx})$ on $[\tau,b]$.

Let $T_*$ be the supremum of times through which this triple
is jointly continuous in $(z,x,t)$.  The short-time family gives
$T_*>a$.  For $k<T_*$ the clamped triple
$\mathcal J_z(x,\min\{t,k\})$ is continuous through $T_*$.
The common time modulus makes these triples converge uniformly
to $\mathcal J_z(x,t)$ as $k\uparrow T_*$, hence continuity
holds through $T_*$ itself.  If $T_*<b$, apply local continuous
dependence to the compact family at $T_*$.  Uniqueness identifies
the new local family with the already chosen individual solutions,
contradicting the supremum.  Therefore $T_*=b$.

Projection preserves continuity of values and first jets, so
$(t,z)\mapsto\pi_Mc_z(\cdot,t)$ is continuous into $\Lambda^1M$.
The initial projected loop is exactly $\widetilde\Gamma(z)$, and the time
path is a homotopy from it to each later loop; nullhomotopy follows.
Apply Proposition~\ref{prop:ramps-length-curvature} to the initial polygon
bounds.  For $0<\rho<1$ these bounds are independent of $\rho$, giving
the displayed exponential estimates and the common constants of
Theorem~\ref{thm:ramps-uniform-jets}.
\end{proof}

The construction applies to every supplied approximation satisfying the
listed conditions at any positive tolerance.  It preserves its count,
polygons and initial family exactly.  The frequently used restriction
$0<\zeta<1$ does not require a new approximation theorem.
Continuity holds for each fixed circumference; joint continuity in $\rho$
is not asserted.

This is the exact interface with the annular-comparison chapter: it receives
the same supplied approximation record, the positive-degree product ramps,
the projected null family, and the circumference-uniform interior estimates.
It must not replace the family by an independently chosen approximation or
identify a product degree with the base-loop homotopy class.  The filling-width
chapter supplies the initial null family, exactly parametrized
near-minimizing disks and uniformly small collars; the next
chapter evolves the corresponding annuli and transfers their area estimates.

\lean{PoincareMT.m63RampEstimates_from_predecessors}
